\documentclass[10pt]{article}
\usepackage[T1]{fontenc}
\usepackage[utf8]{inputenc}
\usepackage[french]{babel}
\usepackage{amsmath,amssymb,mathrsfs,mathpazo}
\usepackage{geometry}
\geometry{a4paper,margin=2cm}
\usepackage{graphicx,array,tabularx,tabularray,colortbl,multirow,multicol,booktabs,enumitem,pifont,ulem,fancybox,diagbox,cancel}
\usepackage{pstricks,pstricks-add,pst-plot,pst-tree,pst-node,pst-text,pst-func,pst-math,pst-bspline,pst-3dplot,pst-eucl,pst-solides3d}
\let\pscalculate\undefined
\usepackage{pst-bezier}
\usepackage{tikz,pgfplots,tkz-tab}
\usetikzlibrary{arrows,arrows.meta,patterns,calc,positioning,trees,decorations.markings,shapes.geometric,shapes.misc,shapes.arrows,intersections,decorations.pathmorphing,decorations.pathreplacing,backgrounds,shadows,babel}
\pgfplotsset{compat=1.18}
\usepgfplotslibrary{fillbetween}
\usepackage[np]{numprint}
\usepackage{listings,makecell,pas-tableur,twoopt}
\let\barre\undefined
\usepackage[pstricks]{tablvar}
\definecolor{darkgreen}{RGB}{0,100,0}
\providecommand{\red}{\color{red}}
\providecommand{\blue}{\color{blue}}
\providecommand{\green}{\color{green}}
\providecommand{\black}{\color{black}}
\providecommand{\white}{\color{white}}
\providecommand{\gray}{\color{gray}}
\providecommand{\cyan}{\color{cyan}}
\providecommand{\magenta}{\color{magenta}}
\providecommand{\yellow}{\color{yellow}}
\providecommand{\violet}{\color{violet}}
\providecommand{\orange}{\color{orange}}
\providecommand{\euro}{\texteuro}
\providecommand{\EUR}{\texteuro}
\providecommand{\ieme}{\textsuperscript{e}}
\providecommand{\degres}{\ensuremath{^\circ}}
\providecommand{\degree}{\ensuremath{^\circ}}
\providecommand{\vv}[1]{\overrightarrow{#1}}
\providecommand{\Coord}[1]{#1}
\providecommand{\decosix}{}
\providecommand{\starredbullet}{\ensuremath{\bullet}}
\pagestyle{empty}
\setlength{\parindent}{0pt}
\usepackage{adjustbox}\DeclareUnicodeCharacter{2236}{:}\DeclareUnicodeCharacter{2190}{\ensuremath{\leftarrow}}
\providecommand{\R}{\mathbb{R}}
\providecommand{\N}{\mathbb{N}}
\providecommand{\Z}{\mathbb{Z}}
\providecommand{\Q}{\mathbb{Q}}
\providecommand{\D}{\mathbb{D}}
\providecommand{\C}{\mathbb{C}}
\providecommand{\e}{{\mathrm{e}}}
\providecommand{\vect}[1]{\overrightarrow{#1}}
\providecommand{\vectt}[1]{\overrightarrow{\mathrm{#1}}}
\providecommand{\vvt}[1]{\overrightarrow{\mathrm{#1}}}
\providecommand{\np}[2][]{{#2~#1}}
\providecommand{\vv}[1]{{\overrightarrow{#1}}}
\renewcommand{\d}{{\mathrm{d}}}
\providecommand{\degres}{{^\circ}}
\providecommand{\degre}{{^\circ}}
\providecommand{\blue}{\color{blue}}
\providecommand{\red}{\color{red}}
\providecommand{\enskip}{\quad}
\providecommand{\llbracket}{\lbrack\!\lbrack}
\providecommand{\rrbracket}{\rbrack\!\rbrack}
\providecommand{\euro}{{\text{€}}}
\providecommand{\no}{{\text{n°}}}
\providecommand{\interval}[2]{{[#1;#2]}}
\providecommand{\intervFO}[2]{{[#1;#2[}}
\providecommand{\dotfill}{\ldots}

\providecommand{\ds}{\displaystyle}
\providecommand{\textding}[1]{#1}
\providecommand{\Oij}{\(\mathrm{O};\vec{\imath},\vec{\jmath}\)}
\providecommand{\Oijk}{\(\mathrm{O};\vec{\imath},\vec{\jmath},\vec{k}\)}
\providecommand{\Ouv}{\(\mathrm{O};\vec{u},\vec{v}\)}
\providecommand{\intervFF}[2]{\ensuremath{[#1;#2]}}
\renewcommand{\euro}{\ifmmode\text{\texteuro}\else\texteuro\fi}
\DeclareUnicodeCharacter{00B0}{\ifmmode^{\circ}\else\textdegree\fi}
\DeclareUnicodeCharacter{20AC}{\ifmmode\text{\texteuro}\else\texteuro\fi}
% Énoncés : Éducation nationale. Transcription et corrigés : APMEP, Denis Vergès et contributeurs. Découpage et classement : CoopMaths, Sébastien Lozano, Éric Elter et contributeurs. Adaptation de la présentation : Maths Explorer.
% CC BY-SA : https://coopmaths.fr/www/about/
% https://coopmaths.fr/alea/static/bac/2024/tex/bac_2024_05_sujet1_ameriquenord_3.tex
\begin{document}
\noindent\textit{APMEP / CoopMaths. CC BY-SA. Présentation adaptée par Maths Explorer.}\par\medskip
\medskip

Le but de cet exercice est d'étudier la fonction $f$ définie sur l'intervalle $]0~;~+\infty[$ par : 

\[f(x) = x\ln \left(x^2\right) -\dfrac 1x.\]

\smallskip

\textbf{Partie A : lectures graphiques}

\medskip

On a tracé ci-dessous la courbe représentative $(\mathcal{C}_f)$ de la fonction $f$, ainsi que la droite $(T)$, tangente à la courbe $(\mathcal{C}_f)$ au point A de coordonnées $(1~;~-1)$.

Cette tangente passe également par le point B$(0~;~-4)$.

\begin{center}
\psset{xunit=1.5cm,yunit=0.6cm,arrowsize=2pt 3}
\begin{pspicture*}(-0.4,-7.2)(4,7.1)
\multido{\n=1+1}{4}{\psline[linewidth=0.2pt](\n,-7)(\n,7)}
\multido{\n=-7+1}{15}{\psline[linewidth=0.2pt](0,\n,)(4,\n)}
\psaxes[linewidth=1.25pt,labelFontSize=\scriptstyle,showorigin=false]{->}(0,0)(-0.4,-7)(4,7)
\psplot[plotpoints=2000,linewidth=1.25pt,linecolor=red]{0.1}{4}{x dup mul ln x mul 1 x div sub}
\psplotTangent{1}{3}{x dup mul ln x mul 1 x div sub}
\uput[l](3,6){\red $\mathcal{C}_f$}\uput[dr](3,5){$T$}
\uput[dr](1,-1){A}\uput[dr](0,-4){B}
\end{pspicture*}
\end{center}

\begin{enumerate}
\item Lire graphiquement $f'(1)$ et donner l'équation réduite de la tangente $(T)$.
\item Donner les intervalles sur lesquels la fonction $f$ semble convexe ou concave.

Que semble représenter le point A pour la courbe $(\mathcal{C}_f)$ ?
\end{enumerate}

\bigskip

\textbf{Partie B : étude analytique}

\medskip

\begin{enumerate}
\item Déterminer, en justifiant, la limite de $f$ en $+\infty$, puis sa limite en $0$.
\item On admet que la fonction $f$ est deux fois dérivable sur l'intervalle $]0~;~+\infty[$.
	\begin{enumerate}
		\item Déterminer $f'(x)$ pour $x$ appartenant à l'intervalle $]0~;~+\infty[$.
		\item Montrer que pour tout $x$ appartenant à l'intervalle $]0~;~+\infty[$,

\[f''(x) = \dfrac{2(x + 1)(x - 1)}{x^3}.\]
	\end{enumerate}
\item 
	\begin{enumerate}
		\item Étudier la convexité de la fonction $f$ sur l'intervalle $]0~;~+\infty[$.
		\item Étudier les variations de la fonction $f'$, puis le signe de $f'(x)$ pour $x$ appartenant à l'intervalle $]0~;~+\infty[$.
		
En déduire le sens de variation de la fonction $f$ sur l'intervalle $]0~;~+\infty[$.
	\end{enumerate}
\item
	\begin{enumerate}
		\item Montrer que l'équation $f(x) = 0$ admet une unique solution $\alpha$ sur l'intervalle $]0~;~+\infty[$.
		\item Donner la valeur arrondie au centième de $\alpha$ et montrer que $\alpha$ vérifie :
		
\[\alpha^2 = \text{exp} \left(\dfrac{1}{\alpha^2}\right).\]
	\end{enumerate}
\end{enumerate}

\bigskip

\end{document}
